\subsection{Visual servoing and guidance for docking}

Visual servo control divides into two schemes. Image-based visual servoing (IBVS) regulates the feature error directly in the image plane, while position-based visual servoing (PBVS) reconstructs a metric relative pose and regulates it in Cartesian space \cite{chaumette2006visual}. In this work, the perception layer already produces a full 6-DOF pose from fiducial markers of known geometry using the standardised marker PnP algorithm \cite{garrido2014automatic}. Therefore, position-based servoing adds no extra cost, and it keeps the docking geometry expressible in physical units, which is why it is adopted here.

However, the servoing literature assumes that the target pose is fixed. For the moving docks of Section~\ref{sec:lit-dock-motion}, the servo target itself oscillates, and regulating onto it raises the estimation question taken up next.
